% Source problem: https://www.ap.t.u-tokyo.ac.jp/themes/ap_t_u-tokyo/graduate/pdf/h30ap_exam.pdf
\documentclass[11pt,a4paper]{article}
\usepackage[margin=24mm]{geometry}
\usepackage{fontspec}
\usepackage{xeCJK}
\setCJKmainfont{HaranoAjiMincho-Regular.otf}
\usepackage{amsmath,amssymb}
\setlength{\parindent}{1em}
\setlength{\parskip}{0.35em}
\title{物理工学専攻 院試解答集\\2018年度 物理学II 第2問}
\author{政岡凜太郎}
\date{}
\begin{document}
\maketitle

\section*{[1.1] 配置数とエントロピー}
$N$個の分子を$M$個のセルのうち異なる$N$個に置くので、配置数は
\begin{equation}
 W=\binom{M}{N}=\frac{M!}{N!(M-N)!}.
\end{equation}
Stirling近似を使うと
\begin{align}
 S&=k_B\ln W\\
 &=k_B\{M\ln M-N\ln N-(M-N)\ln(M-N)\}\\
 &=-Mk_B\left[\phi\ln\phi+(1-\phi)\ln(1-\phi)\right],
 \qquad \phi=\frac NM.
\end{align}

\section*{[1.2] 圧力と理想気体との比較}
相互作用がないため$F_0=-TS$である。$v$を一定に保って$V$を変えると、$M=V/v$かつ$\phi=vN/V$も変わる。$f(\phi)=\phi\ln\phi+(1-\phi)\ln(1-\phi)$と置けば
\begin{equation}
 F_0=\frac{k_BT V}{v}f(\phi),\qquad
 \left(\frac{\partial\phi}{\partial V}\right)_{N}=-\frac{\phi}{V}.
\end{equation}
従って
\begin{align}
 p_0=-\left(\frac{\partial F_0}{\partial V}\right)_{T,N}
 &=-\frac{k_BT}{v}\left[f(\phi)-\phi f'(\phi)\right]\\
 &=-\frac{k_BT}{v}\ln(1-\phi)
 =\boxed{-\frac{k_BT}{v}\ln\left(1-\frac{vN}{V}\right)}.
\end{align}
$0<\phi<1$では対数の展開から
\begin{equation}
 p_0=\frac{k_BT}{v}\left(\phi+\frac{\phi^2}{2}+\cdots\right)
 =\frac{k_BNT}{V}+\frac{k_BTvN^2}{2V^2}+\cdots.
\end{equation}
希薄極限の一次項は$p_{\rm id}=k_BNT/V$に一致する。有限密度ではセルの二重占有が禁じられているため、自由に使える空きセルが減って圧力は理想気体より大きくなる。

\section*{[2.1] 平均場化学ポテンシャル}
[1.1]のエントロピーと与えられた平均場エネルギーより
\begin{equation}
 F(\phi)=-\frac12 Mz\alpha\phi^2+Mk_BT\left[\phi\ln\phi+(1-\phi)\ln(1-\phi)\right].
\end{equation}
$N=M\phi$なので、$M$を固定したまま$\mu=(\partial F/\partial N)_T=M^{-1}(\partial F/\partial\phi)_T$を計算すると
\begin{equation}
 \boxed{\mu(\phi)=-\alpha z\phi+k_BT\left[\ln\phi-\ln(1-\phi)\right]}.
\end{equation}
第一項は近接分子間の引力による寄与、対数項は占有数の制約によるエントロピーの寄与である。

\section*{[2.2] 化学ポテンシャルが単調増加する温度}
\begin{equation}
 \mu'(\phi)=-\alpha z+\frac{k_BT}{\phi(1-\phi)}.
\end{equation}
$0<\phi<1$における$\phi(1-\phi)$の最大値は$1/4$である。よって$\mu'(\phi)\ge0$が全域で成り立つ条件は
\begin{equation}
 \boxed{T\ge T_c\equiv\frac{\alpha z}{4k_B}}.
\end{equation}
$T>T_c$では導関数は常に正であり、$T=T_c$では$\phi=1/2$の一点だけでゼロ、それ以外では正である。従って臨界温度でも$\mu(\phi)$は区間全体で厳密に増加する。

\section*{[2.3] 三つの密度と相分離}
$T<T_c$では$\mu'(\phi)=0$となる密度が二つあり、その間で$\mu$は減少する。このS字形の曲線では、極小値と極大値の間のある化学ポテンシャルに対し、同じ$\mu$を与える密度が三つ現れる。同じ大きさ$M$の二領域に粒子を$N+\delta N$と$N-\delta N$だけ分けたとき、均等な配置との差は
\begin{equation}
 F(N+\delta N)+F(N-\delta N)-2F(N)
 =F''(N)(\delta N)^2+O((\delta N)^4)
\end{equation}
である。各$F$は同じ$M,T$で評価する。したがって安定な一様相では$F''(N)=\partial\mu/\partial N\ge0$が必要である。中央の枝では$\partial\mu/\partial N=M^{-1}\mu'(\phi)<0$となり、粒子を偏らせると自由エネルギーが下がるため、この安定性条件に反する。

安定な平衡では、組成に応じて低密度相と高密度相に分かれることがある。$g(\phi)=F(\phi)/M$をセルあたりの自由エネルギーとすれば、共存密度$\phi_g,\phi_l$は
\begin{equation}
 g'(\phi_g)=g'(\phi_l)=\mu_{\rm coex},\qquad
 g(\phi_l)-g(\phi_g)=\mu_{\rm coex}(\phi_l-\phi_g)
\end{equation}
を満たす。第一式は両相の化学ポテンシャルの等値、第二式は二点を結ぶ直線の傾きが共通の接線勾配に等しいことを表す。固定セル体積$v$では圧力は$p=(\phi\mu-g)/v$なので、二式を合わせると圧力も等しい。第二式を積分形に直すと$\int_{\phi_g}^{\phi_l}[\mu(\phi)-\mu_{\rm coex}]\,d\phi=0$となり、Maxwellの等面積則を得る。これが液体と気体の相分離に対応する。三つの解は三相が同時に共存することを意味せず、両端の局所的に安定な枝と中央の不安定枝を表している。
\end{document}
